\documentclass[a4paper,UTF8,zihao=-4,fontset=none]{ctexart}
\usepackage[margin=2.5cm]{geometry}
\setCJKmainfont[Path=fonts/,FontIndex=0,AutoFakeBold=true]{simsun.ttc}
\setCJKsansfont[AutoFakeBold=true]{FandolHei-Regular.otf}
\usepackage{hyperref}
\hypersetup{hidelinks,pdftitle={AI工具使用详情},pdfauthor={}}
\setlength{\parindent}{2em}
\begin{document}
\begin{center}\Large\bfseries AI工具使用详情\end{center}
\section*{1. 工具与记录范围}
使用工具为OpenAI Codex，包含对话式推理、代码编写与执行、网页及文献查询、浏览器编辑等功能。本次论文整理会话基于GPT-6；此前各次开发会话的具体型号未全部保留，本说明不推定其与本次相同。记录范围为本项目现有可核对的模型、程序及论文协作过程，整理日期为2026年9月11日。

\section*{2. 使用目的与环节}
AI辅助讨论测向扇区与半平面交、定位区域直径和最小覆盖圆的区别；协助推导保证再次收信的候选域、最坏相容读数目标及数值上界；辅助实现与调试四问程序，构造本地样本、比较搜索策略并整理结果。论文阶段使用AI核查原题和实现的一致性、查找并核对相关文献、生成图表、撰写和压缩中文表述，以及在Overleaf整理LaTeX格式。

AI参与范围包括模型讨论和程序实现，并不限于语言润色。正式平台成绩必须由实际测试产生，AI没有以自行构造的数据替代正式测试。

\section*{3. 主要提示方式与过程}
提示采用多轮对话，提供题目、原始附件、当前代码目录和已有实验结果，要求说明依据、区分理论保证与有限样本结果，禁止策略读取真实源位置或场景类别。典型提示包括：“模仿数学建模论文的语言，控制篇幅”；“结合之前我们讨论过的模型假设以及你运行代码的实际假设”；“根据当前进度把论文剩余部分都写完，参考文献也加好，原文中引用一下。格式要求见官方的文件”。

论文整合中，先按原题及仓库实现核对四问的假设、算法和结果，再查阅作者稿、出版社及官方文件，编写对应公式、结果表与引用；最后将正文和完整程序附录编译为PDF并检查排版。模型核对、文献核对和图表整理使用了Codex中的并行辅助工作单元，其输出由主会话交叉检查后合并。

\section*{4. 采纳、修改与核验情况}
采纳内容包括经核对的集合定位表述、与当前程序一致的算法说明、现有结果文件中的统计值，以及可验证题录对应的局部方法引用。整理时特别检查了四圆盘收信条件、同地点固定误差、接口读数舍入、方向源阴性裁剪的前提，以及发现16个频道与实际清除16个源的区别；没有将未采用的第三轮速度实验作为正文默认算法。

程序自动核验方面，现有统一报告记录第一至四问分别通过36、27、41、38项测试，并含异目录启动、结果一致性、困难样本回归及资料完整性检查。论文还核对了第二问上界评分与实际后验的共同约束，以及提前停止与完整扫描得到相同评分和选点的结果。上述均属于程序验证或AI辅助复核，不能等同于参赛者已经完成的逐项人工审查。

人工修改方面，现有对话记录显示参赛者持续指定选题范围、建模与验证要求、论文结构和篇幅，并要求与实际实现一致；这些要求已用于筛选和修订内容。现有记录不足以证明参赛者已逐项审查本次新生成的全文与全部代码，因此本说明不作此项已完成的声明。正式提交前，参赛队仍须逐项人工核实采纳的核心推导、代码、统计表、参考文献和AI使用记录，并据实更新本段。

当前证据尚不足以填报第三、四问各三次正式测试成绩。若后续补入成绩，应同步核对案例编码、平台时间和原名加密日志，更新正文、支撑文件清单与本说明。此文件不代替参赛队的最终审查和确认。
\end{document}
